\chapter{Filtracje, adaptacja i czasy zatrzymania}
\label{ch:filtracje}

\section{Informacja rosnąca w czasie}

Warunkowa wartość oczekiwana (rozdział~\ref{ch:condexp})
opisuje to, co można przewidzieć na podstawie ustalonej
informacji.
W procesie losowym informacja jednak rośnie:
po każdym kroku poznajemy nowy wynik. \emph{Filtracja}
porządkuje te kolejne etapy, a warunek adaptacji
gwarantuje, że wartość procesu w chwili $n$
jest mierzalna względem przyjętej informacji
$\mathcal F_n$. To, czy informacja ta zawiera
również pewne dane o przyszłości, zależy od
wyboru filtracji.

\begin{definicja}[Filtracja i proces adaptowany]
Niech $(\Omega,\mathcal F,\Prob)$ będzie przestrzenią
probabilistyczną. \emph{Filtracją w czasie dyskretnym}
nazywamy ciąg $\sigma$-ciał
$(\mathcal F_n)_{n\geq0}$ spełniający
$\mathcal F_n\subseteq\mathcal F_{n+1}\subseteq\mathcal F$
dla każdego $n\geq0$.
Procesem o wartościach w przestrzeni mierzalnej
$(S,\mathcal S)$ nazywamy tutaj ciąg mierzalnych
odwzorowań $X_n:\Omega\to S$.
Proces $(X_n)_{n\geq0}$ jest \emph{adaptowany}
względem $(\mathcal F_n)$, jeśli każde $X_n$
jest $(\mathcal F_n,\mathcal S)$-mierzalne.
\emph{Naturalna filtracja} procesu to
$\mathcal F_n^X=\sigma(X_0,\ldots,X_n)$.
Proces $(H_n)_{n\geq1}$ o wartościach w ustalonej
przestrzeni mierzalnej jest \emph{przewidywalny}
względem $(\mathcal F_n)$, jeśli każde $H_n$
jest mierzalne względem $\mathcal F_{n-1}$.
\end{definicja}

\begin{uwaga}[Najmniejsza filtracja i przewidywalność]
Naturalna filtracja jest filtracją, ponieważ dołączenie
kolejnej zmiennej nie zmniejsza generowanego
$\sigma$-ciała. Jeżeli $X$ jest adaptowany względem
$(\mathcal G_n)$, to dla $k\leq n$ zmienna $X_k$
jest mierzalna względem $\mathcal G_k\subseteq\mathcal G_n$,
zatem $\mathcal F_n^X\subseteq\mathcal G_n$.
W szczególności naturalna filtracja jest najmniejszą
filtracją, względem której $X$ jest adaptowany.
Każdy proces przewidywalny jest adaptowany w chwilach
$n\geq1$, gdyż $\mathcal F_{n-1}\subseteq\mathcal F_n$.
Odwrotna implikacja nie zachodzi, jak pokazuje
następujący przykład z niezależnymi rzutami monetą.
\end{uwaga}

\begin{przyklad}[Historia, dodatkowa informacja i uprzedzenie przyszłości]
Dla niezależnych uczciwych rzutów monetą
$\xi_k\in\{-1,1\}$,
$S_n=\sum_{k=1}^n\xi_k$ i $S_0=0$, naturalna
filtracja sum $\mathcal F_n^S$ jest równa
$\sigma(\xi_1,\ldots,\xi_n)$: sumy są mierzalnymi
funkcjami przyrostów, a z trajektorii $S_0,\ldots,S_n$
odzyskujemy wszystkie przyrosty $\xi_k=S_k-S_{k-1}$.
Dla $n=0$ otrzymujemy trywialne $\sigma$-ciało
$\{\varnothing,\Omega\}$.
Możemy też rozszerzyć informację o dodatkowe
obserwacje $Y_0,\ldots,Y_n$, przyjmując
$\mathcal G_n=\sigma(S_0,\ldots,S_n,Y_0,\ldots,Y_n)$.
W szczególności przyjmijmy
$\mathcal G_n=\sigma(\xi_1,\ldots,\xi_{\max(1,n)})$.
Proces sum pozostaje adaptowany, ale nie jest już
martyngałem względem $(\mathcal G_n)$, gdyż
$\E[S_1\mid\mathcal G_0]=\xi_1\ne0=S_0$ p.n.
W naturalnej filtracji jest martyngałem, ponieważ
niezależny od wcześniejszych rzutów przyrost ma
średnią zero (dowód podamy również niżej). Martyngałowość
zależy więc \emph{od filtracji}, a nie tylko od
wspólnego rozkładu procesu.
\end{przyklad}

\begin{przyklad}[Strategia przewidywalna]
Jeżeli $H_n$ oznacza liczbę jednostek aktywa,
którą wybieramy \emph{przed} zaobserwowaniem
zmiany ceny między $n-1$ i $n$, to powinna być
$\mathcal F_{n-1}$-mierzalna. Strategia
$H_n=\operatorname{sgn}(S_n-S_{n-1})$ używałaby
już wyniku kroku $n$. W naturalnej filtracji
niezależnych uczciwych rzutów z poprzedniego przykładu
jest to $H_n=\xi_n$. Gdyby $H_n$ było
$\mathcal F_{n-1}$-mierzalne, zdarzenie $\{\xi_n=1\}$
byłoby niezależne od siebie, a jego prawdopodobieństwo
$1/2$ musiałoby spełniać $1/2=(1/2)^2$.
Otrzymujemy sprzeczność. Jest to więc proces adaptowany,
ale nieprzewidywalny względem tej filtracji.
\end{przyklad}


\begin{przyklad}[Filtracja rozpisana na elementarnych zdarzeniach]
Wykonujemy dwa niezależne uczciwe rzuty monetą.
Przestrzeń wyników jest
$\Omega=\{OO,OR,RO,RR\}$ z $\mathcal F=2^\Omega$,
gdzie pierwsza litera
opisuje pierwszy rzut. W chwili $0$ nic nie wiemy:
$\mathcal F_0=\{\varnothing,\Omega\}$. Po pierwszym
rzucie rozróżniamy atomy $\{OO,OR\}$ i $\{RO,RR\}$,
a więc
\[
\mathcal F_1=\{\varnothing,\Omega,
\{OO,OR\},\{RO,RR\}\},\qquad
\mathcal F_2=2^\Omega.
\]
Dla $n\geq2$ przyjmujemy $\mathcal F_n=\mathcal F_2$.
Niech $Y=\chi_{\{OO\}}$ będzie funkcją
charakterystyczną zdarzenia, że oba rzuty dały orła.
Każdy wynik ma prawdopodobieństwo $1/4$.
Obliczając średnie na atomach odpowiednich
$\sigma$-ciał (rozdział~\ref{ch:condexp}), otrzymujemy
$\E[Y\mid\mathcal F_0]=1/4$,
$\E[Y\mid\mathcal F_1]=(1/2)\chi_{\{OO,OR\}}$
oraz $\E[Y\mid\mathcal F_2]=Y$.
Pierwsza prognoza jest stała, druga wykorzystuje wynik
pierwszego rzutu, a trzecia jest już dokładną
obserwacją. Pokazuje to, dlaczego filtracja
to matematyczny model \emph{narastania informacji}
(por. rysunek~\ref{fig:rp1-13-filtration}).
\end{przyklad}

\begin{figure}[htbp]
\centering
\begin{tikzpicture}[>=Stealth,
  atom/.style={draw,rounded corners=3pt,minimum height=8mm,
  align=center,font=\small,fill=blue!6,draw=blue!65!black}]
\node[atom,minimum width=7cm] (a) at (0,0) {$\{OO,OR,RO,RR\}$};
\node[atom,minimum width=3.1cm] (b) at (-1.9,-1.5) {$\{OO,OR\}$};
\node[atom,minimum width=3.1cm] (c) at (1.9,-1.5) {$\{RO,RR\}$};
\foreach \x/\s in {-2.8/OO,-1/OR,1/RO,2.8/RR}
  \node[atom,minimum width=1.5cm] (\s) at (\x,-3) {$\{\s\}$};
\draw[->,thick] (a.south)--(b.north);
\draw[->,thick] (a.south)--(c.north);
\draw[->] (b.south)--(OO.north);
\draw[->] (b.south)--(OR.north);
\draw[->] (c.south)--(RO.north);
\draw[->] (c.south)--(RR.north);
\node[anchor=east] at (-4,0) {$\mathcal F_0$};
\node[anchor=east] at (-4,-1.5) {$\mathcal F_1$};
\node[anchor=east] at (-4,-3) {$\mathcal F_2$};
\end{tikzpicture}
\caption{Uszczegóławianie podziału przestrzeni wyników
przez kolejne rzuty. W każdym wierszu pokazano
atomy danego $\sigma$-ciała; jego zdarzeniami
są wszystkie sumy tych atomów.}
\label{fig:rp1-13-filtration}
\end{figure}


\section{Czasy zatrzymania i informacja przy zatrzymaniu}

\begin{definicja}[Czas zatrzymania]
Zmienna losowa $\tau:\Omega\to\N_0\cup\{\infty\}$
jest \emph{czasem zatrzymania} względem filtracji
$(\mathcal F_n)$, jeśli $\{\tau\leq n\}\in\mathcal F_n$
dla każdego $n\geq0$. Równoważnie
$\{\tau=n\}\in\mathcal F_n$ dla każdego $n\geq0$:
dla $n=0$ oba zdarzenia są identyczne, dla $n\geq1$
odejmujemy $\{\tau\leq n-1\}$ od $\{\tau\leq n\}$,
a w drugą stronę korzystamy z tożsamości
$\{\tau\leq n\}=\bigcup_{k=0}^n\{\tau=k\}$.
Warunki te zapewniają również $\mathcal F$-mierzalność
$\tau$, ponieważ $\{\tau=\infty\}$ jest dopełnieniem
przeliczalnej sumy zdarzeń $\{\tau=n\}$.
\end{definicja}

\begin{przyklad}[Czas wejścia do zbioru]
Dla procesu adaptowanego $(X_n)$ o wartościach
w przestrzeni mierzalnej $(S,\mathcal S)$ i
$A\in\mathcal S$ czas
$\tau_A=\inf\{n\geq0:X_n\in A\}$, z konwencją
$\inf\varnothing=\infty$, jest czasem zatrzymania, bo
$\{\tau_A\leq n\}=\bigcup_{k=0}^n\{X_k\in A\}
\in\mathcal F_n$.
\end{przyklad}

\begin{przyklad}[Ostatnia wizyta na ogół nie jest czasem zatrzymania]
Dla niezależnych uczciwych rzutów monetą z wartościami
$\xi_k\in\{-1,1\}$, ich naturalnej filtracji
$\mathcal F_n=\sigma(\xi_1,\ldots,\xi_n)$
i ustalonego $N\geq2$
rozważmy $\rho=\max\{k\in\{1,\ldots,N\}:\xi_k=1\}$,
a gdy nie ma takiego $k$, przyjmijmy $\rho=0$.
Zdarzenie $\{\rho\leq N-1\}=\{\xi_N=-1\}$
nie należy do $\sigma(\xi_1,\ldots,\xi_{N-1})$,
ponieważ ostatni rzut jest niezależny od wcześniejszych
i $\Prob(\xi_N=-1)=1/2$. Gdyby to zdarzenie
należało do wcześniejszego $\sigma$-ciała, z niezależności
wynikałoby, że jego prawdopodobieństwo jest równe
własnemu kwadratowi, co jest niemożliwe.
Dlatego $\rho$ nie jest czasem zatrzymania. Aby rozpoznać \emph{ostatni} sukces,
trzeba znać jeszcze przyszłość.
\end{przyklad}

\begin{definicja}[Zatrzymana informacja]
Dla czasu zatrzymania $\tau$ definiujemy
\[
\mathcal F_\tau=
\{A\in\mathcal F: A\cap\{\tau\leq n\}\in\mathcal F_n
\text{ dla każdego }n\geq0\}.
\]
Jest to $\sigma$-ciało: zawiera $\Omega$, a dla
$A\in\mathcal F_\tau$ mamy
$A^c\cap\{\tau\leq n\}=\{\tau\leq n\}\setminus
(A\cap\{\tau\leq n\})\in\mathcal F_n$.
Przecięcie z $\{\tau\leq n\}$ rozdziela się względem
przeliczalnej sumy zbiorów, więc $\mathcal F_\tau$
jest także zamknięte na takie sumy.
Na $\{\tau<\infty\}$ opisuje ono informację dostępną
w chwili zatrzymania.
Dla mierzalnego procesu rzeczywistego przyjmujemy
\[
X_\tau(\omega)=
\begin{cases}
X_{\tau(\omega)}(\omega),&\tau(\omega)<\infty,\\
0,&\tau(\omega)=\infty.
\end{cases}
\]
Równoważnie $X_\tau=\sum_{k\geq0}X_k\chi_{\{\tau=k\}}$,
gdzie w każdym punkcie co najwyżej jeden składnik
jest niezerowy. Proces zatrzymany definiujemy przez
$X^\tau_n=X_{n\wedge\tau}$; tutaj indeks jest zawsze
skończony.
\end{definicja}

\begin{lemat}[Podstawowe własności zatrzymania]
Niech $\tau,\sigma$ będą czasami zatrzymania
względem $(\mathcal F_n)$. Wtedy:
\begin{enumerate}
\item $\tau\wedge\sigma$ oraz $\tau\vee\sigma$ są czasami zatrzymania;
\item jeśli $\sigma\leq\tau$ punktowo, to
  $\mathcal F_\sigma\subseteq\mathcal F_\tau$;
\item $\tau$ jest $\mathcal F_\tau$-mierzalny;
\item jeśli $X$ jest rzeczywistym procesem adaptowanym,
  to $X_\tau$ (z umowną wartością zero na
  $\{\tau=\infty\}$) jest $\mathcal F_\tau$-mierzalna,
  a proces zatrzymany $(X_{n\wedge\tau})$ jest adaptowany.
\end{enumerate}
\end{lemat}
\begin{proof}
Dla minimum mamy
$\{\tau\wedge\sigma\leq n\}=
\{\tau\leq n\}\cup\{\sigma\leq n\}$,
a dla maksimum zastępujemy sumę zbiorów
przecięciem. Oba należą do $\mathcal F_n$.
Jeśli $A\in\mathcal F_\sigma$, to przy
$\sigma\leq\tau$ zachodzi równość
\[
A\cap\{\tau\leq n\}
=(A\cap\{\sigma\leq n\})\cap\{\tau\leq n\},
\]
wynikająca z $\sigma\leq\tau$, a prawa strona
należy do $\mathcal F_n$.
Dla (3) wystarczy zauważyć, że
$\{\tau\leq m\}\cap\{\tau\leq n\}
=\{\tau\leq\min(m,n)\}\in\mathcal F_n$,
więc $\{\tau\leq m\}\in\mathcal F_\tau$ dla każdego $m$.
Z tych zdarzeń przez różnice i przeliczalne
sumy oraz dopełnienie otrzymujemy przeciwobraz
każdego podzbioru $\N_0\cup\{\infty\}$.
Dla (4) niech $B\subseteq\R$ będzie borelowski.
Na zdarzeniu $\{\tau\leq n\}$ przeciwobraz
$\{X_\tau\in B\}$ jest sumą skończoną
$\bigcup_{k=0}^n\{\tau=k\}\cap\{X_k\in B\}$,
należącą do $\mathcal F_n$. Stąd mierzalność
względem $\mathcal F_\tau$, gdyż $X_\tau$
jest także $\mathcal F$-mierzalna jako punktowa
granica sum częściowych wskazanych w definicji.
Ponieważ $n\wedge\tau$ jest czasem zatrzymania
nie większym niż $n$, z (2) mamy
$\mathcal F_{n\wedge\tau}\subseteq\mathcal F_n$.
Tutaj $\mathcal F_n$ jest zarazem $\sigma$-ciałem
zatrzymanym w deterministycznym czasie $n$:
dla $k<n$ warunek definicyjny jest pusty,
dla $k=n$ wymaga dokładnie $A\in\mathcal F_n$,
a dla $k>n$ jest wtedy spełniony na mocy
$\mathcal F_n\subseteq\mathcal F_k$.
Stąd $X_{n\wedge\tau}$ jest $\mathcal F_n$-mierzalna.
\end{proof}

\begin{uwaga}[Dwie równoważne definicje $\mathcal F_\tau$]
Dla każdego $A\in\mathcal F$ warunki
$A\cap\{\tau\leq n\}\in\mathcal F_n$ dla wszystkich $n$
oraz $A\cap\{\tau=n\}\in\mathcal F_n$ dla wszystkich $n$
są równoważne. Dla $n=0$ warunki są identyczne,
a dla $n\geq1$ korzystamy z różnicy
$\{\tau=n\}=\{\tau\leq n\}\setminus\{\tau\leq n-1\}$
i z inkluzji $\mathcal F_{n-1}\subseteq\mathcal F_n$.
W drugą stronę korzystamy z sumy
$\{\tau\leq n\}=\bigcup_{k=0}^n\{\tau=k\}$.
Żadnego dodatkowego warunku na $\{\tau=\infty\}$
nie trzeba dopisywać do tej równoważności.
W przyjętej definicji, gdy $\tau\equiv\infty$,
mamy nawet $\mathcal F_\tau=\mathcal F$:
proces nie został zatrzymany w żadnej skończonej chwili,
a definicja nie nakłada ograniczeń na zdarzenia
poza częściami $\{\tau\leq n\}$. Jest to konwencja
wygodna dla ogólnych czasów zatrzymania;
przy zatrzymaniu skończonym p.n. dowolne
ograniczenie informacji na $\{\tau=\infty\}$
nie zmienia warunkowania p.n., o ile pozostawia
wszystkie zdarzenia na $\{\tau<\infty\}$.
Istotnie, różnice takich zdarzeń są zawarte
w zbiorze miary zero. Nie należy jednak bez
dodatkowych założeń utożsamiać samych $\sigma$-ciał.
\end{uwaga}


\begin{przyklad}[Jaka jest informacja w chwili losowego zatrzymania?]
Niech $\xi_1,\xi_2\in\{O,R\}$ będą niezależnymi,
uczciwymi rzutami monetą, a
$\tau=1$ przy $\xi_1=O$ oraz $\tau=2$ przy $\xi_1=R$.
Filtracja $\mathcal F_0,\mathcal F_1,\mathcal F_2$
pochodzi z poprzedniego przykładu czteroelementowego.
Jest to czas zatrzymania, ponieważ
$\{\tau\leq0\}=\varnothing$, $\{\tau\leq1\}=\{OO,OR\}$
i $\{\tau\leq2\}=\Omega$.
Z definicji $\mathcal F_\tau$ zbiór $A$ należy
do niego dokładnie wtedy, gdy $A\cap\{OO,OR\}$
jest pusty lub równy $\{OO,OR\}$ (warunek przy
$n=2$ jest automatyczny).
Atomami $\mathcal F_\tau$ są więc
$\{OO,OR\}$, $\{RO\}$ i $\{RR\}$.
Na pierwszym atomie przestajemy obserwować po
pierwszym rzucie i nie rozróżniamy wyników
$OO$ i $OR$. Na dwóch pozostałych atomach
widzieliśmy również drugi rzut.
Dla $Y=\chi_{\{OO\}}$ średnia na atomie
$\{OO,OR\}$ wynosi $1/2$, a na każdym
z dwóch pozostałych atomów jest zerowa. Otrzymujemy
$\E[Y\mid\mathcal F_\tau]
=\frac12\chi_{\{OO,OR\}}$.
Przykład uświadamia, że $\mathcal F_\tau$ nie jest
jednym z góry ustalonym $\mathcal F_n$:
zakres obserwacji zależy od ścieżki, co przedstawia
rysunek~\ref{fig:rp1-13-stopped-tree}.
\end{przyklad}

\begin{figure}[htbp]
\centering
\begin{tikzpicture}[>=Stealth,
  history/.style={draw,rounded corners=3pt,minimum height=9mm,
  minimum width=1.7cm,align=center,font=\small},
  stop/.style={history,fill=green!12,draw=green!45!black,very thick}]
\node[history] (root) at (0,0) {$\Omega$};
\node[stop,minimum width=3.1cm] (o) at (3,1.3) {$\{OO,OR\}$};
\node[history] (r) at (3,-1.3) {$\{RO,RR\}$};
\node[stop] (ro) at (6,-.55) {$\{RO\}$};
\node[stop] (rr) at (6,-2.05) {$\{RR\}$};
\draw[->,thick] (root)--node[above,sloped] {$O$} (o);
\draw[->,thick] (root)--node[below,sloped] {$R$} (r);
\draw[->,thick] (r)--node[above,sloped] {$O$} (ro);
\draw[->,thick] (r)--node[below,sloped] {$R$} (rr);
\node[font=\small,above=2mm of o] {$\tau=1$};
\node[font=\small,right=2mm of ro] {$\tau=2$};
\node[font=\small,right=2mm of rr] {$\tau=2$};
\node at (0,2.5) {$n=0$};
\node at (3,2.5) {$n=1$};
\node at (6,2.5) {$n=2$};
\end{tikzpicture}
\caption{Drzewo obserwacji zatrzymanej: po pierwszym orle
nie wykonujemy dalszej obserwacji, a po reszce obserwujemy
jeszcze drugi rzut. Wyróżnione liście są dokładnie
atomami $\mathcal F_\tau$.}
\label{fig:rp1-13-stopped-tree}
\end{figure}


\section{Przewidywanie przy losowym czasie}

\begin{twierdzenie}[Warunkowanie po zatrzymaniu]
\label{thm:rp1-13-stopped-condexp}
Niech $\tau$ będzie czasem zatrzymania takim,
że $\Prob(\tau<\infty)=1$, niekoniecznie
ograniczonym. Dla $X\in L^1(\Prob)$ oraz
$Z_n=\E[X\mid\mathcal F_n]$, $n\geq0$, zachodzi
\[
\E[X\mid\mathcal F_\tau]
=Z_\tau=\sum_{n\geq0}\chi_{\{\tau=n\}}Z_n
\quad\text{p.n.}
\]
W szczególności, jeśli $\tau\leq N$ p.n. dla
pewnego $N\in\N_0$, sumę można zakończyć na $N$.
\end{twierdzenie}
\begin{proof}
Wybierzmy $\mathcal F_n$-mierzalne rzeczywiste
wersje $Z_n$. Prawa strona jest
$\mathcal F_\tau$-mierzalna na mocy lematu
o zatrzymywaniu zastosowanego do adaptowanego
procesu $(Z_n)$, przy umownej wartości zero
na $\{\tau=\infty\}$.
Ponieważ $|Z_n|\leq\E[|X|\mid\mathcal F_n]$ p.n.
na mocy warunkowej nierówności Jensena
(rozdział~\ref{ch:condexp}), a
$\{\tau=n\}\in\mathcal F_n$, mamy
\[
\E[|Z_n|\chi_{\{\tau=n\}}]
\leq\E[|X|\chi_{\{\tau=n\}}].
\]
Sumując nieujemne wyrazy i stosując twierdzenie
o zbieżności monotonicznej, otrzymujemy
\[
\E|Z_\tau|
=\sum_{n\geq0}\E[|Z_n|\chi_{\{\tau=n\}}]
\leq\E|X|<\infty.
\]
Dla dowolnego $A\in\mathcal F_\tau$ mamy
$A\cap\{\tau=n\}\in\mathcal F_n$. Dlatego
\begin{align*}
\E[\chi_A Z_\tau]
&=\sum_{n\geq0}\E[\chi_{A\cap\{\tau=n\}}Z_n]\\
&=\sum_{n\geq0}\E[\chi_{A\cap\{\tau=n\}}X]
=\E[\chi_A X].
\end{align*}
Oba szeregi całek są absolutnie zbieżne
na mocy powyższego oszacowania oraz całkowalności
$X$, a ostatnia równość wykorzystuje
$\Prob(\tau<\infty)=1$. Otrzymana tożsamość
dla wszystkich $A\in\mathcal F_\tau$ wraz
z mierzalnością i całkowalnością $Z_\tau$
dowodzi tezy.
\end{proof}

\begin{uwaga}
Powyższe twierdzenie dotyczy procesu
$Z_n=\E[X\mid\mathcal F_n]$ pochodzącego
od jednej całkowalnej zmiennej $X$.
Nie wynika z niego twierdzenie o dowolnym
martyngale zatrzymanym w prawie na pewno
skończonym czasie. Odpowiednie dodatkowe warunki
dla martyngałów omówimy poniżej oraz w rozdziale~\ref{ch:nadmartyngaly}.
\end{uwaga}

\section{Martyngał jako konsekwencja warunkowania}

\begin{definicja}[Martyngał, podmartyngał, nadmartyngał]
Niech $(\mathcal F_n)$ będzie filtracją i
$(M_n)$ rzeczywistym procesem adaptowanym
spełniającym $\E|M_n|<\infty$ dla każdego $n$.
Nazywamy go \emph{martyngałem}, jeśli
$\E[M_{n+1}\mid\mathcal F_n]=M_n$ p.n.
dla każdego $n\geq0$;
\emph{podmartyngałem}, jeśli warunkowa wartość
oczekiwana jest większa lub równa $M_n$;
\emph{nadmartyngałem}, jeśli jest mniejsza
lub równa $M_n$. Wszystkie nierówności
rozumiemy prawie na pewno.
\end{definicja}

W tym tomie słowo \emph{martyngał}
bez dopowiedzenia dotyczy czasu dyskretnego
$\N_0$. W czasie ciągłym dochodzą zagadnienia
regularności trajektorii i modyfikacji procesu.
Z mierzalności $M_n$ i liniowości warunkowej
wartości oczekiwanej wynika równoważny warunek
$\E[M_{n+1}-M_n\mid\mathcal F_n]=0$ p.n.;
dla podmartyngałów warunkowa średnia przyrostu
jest nieujemna, a dla nadmartyngałów niedodatnia.

\begin{przyklad}[Ruch bez dryfu i z dryfem]
Niech $(\xi_k)$ będą niezależne i całkowalne,
$\mathcal F_n=\sigma(\xi_1,\ldots,\xi_n)$,
$S_n=\sum_{k=1}^n\xi_k$ i $S_0=0$.
Każde $S_n$ jest $\mathcal F_n$-mierzalne
i całkowalne, gdyż $\E|S_n|\leq\sum_{k=1}^n\E|\xi_k|$.
Z niezależności $\xi_{n+1}$ od $\mathcal F_n$ wynika
$\E[\xi_{n+1}\mid\mathcal F_n]=\E\xi_{n+1}$.
Jeśli $\E\xi_k=0$ dla każdego $k$, to $(S_n)$
jest martyngałem, ponieważ
$\E[S_{n+1}\mid\mathcal F_n]
=S_n+\E\xi_{n+1}=S_n$.
Jeśli $\E\xi_k=m$ dla wszystkich $k$, to
dla $m>0$ proces $(S_n)$ jest podmartyngałem,
a dla $m<0$ nadmartyngałem. Bez założenia
o wspólnej średniej proces
$S_n-\sum_{k=1}^n\E\xi_k$ jest martyngałem:
jego przyrost ma warunkową średnią zero.
Dla wspólnej średniej $m$ otrzymujemy więc
martyngał $(S_n-nm)$. Zerowa warunkowa średnia
przyrostu nie oznacza niezależności przyrostu
od przeszłości.
\end{przyklad}

\begin{twierdzenie}[Kryterium wszystkich przyszłych chwil]
Jeśli $(M_n)$ jest procesem całkowalnym i
adaptowanym, to jest martyngałem wtedy i tylko wtedy,
gdy dla wszystkich liczb $0\leq m\leq n$
$\E[M_n\mid\mathcal F_m]=M_m$ p.n.
Analogicznie podmartyngały i nadmartyngały
spełniają odpowiednią nierówność dla $m\leq n$.
\end{twierdzenie}
\begin{proof}
Implikacja odwrotna wynika z wyboru $n=m+1$.
Ustalmy $m$ i dowiedźmy implikacji wprost
indukcyjnie względem $n\geq m$.
Dla $n=m$ korzystamy z adaptacji:
$\E[M_m\mid\mathcal F_m]=M_m$ p.n.
Dla $n>m$ własność wieży daje
$\E[M_n\mid\mathcal F_m]
=\E[\E[M_n\mid\mathcal F_{n-1}]\mid\mathcal F_m]$.
Dla martyngału równa się to
$\E[M_{n-1}\mid\mathcal F_m]=M_m$.
Dla podmartyngału używamy monotoniczności
warunkowej wartości oczekiwanej oraz założenia
indukcyjnego, otrzymując
$\E[M_n\mid\mathcal F_m]\geq
\E[M_{n-1}\mid\mathcal F_m]\geq M_m$.
Dla nadmartyngału nierówności odwracamy.
\end{proof}

\section{Zatrzymywanie procesu: podstawowa konstrukcja}

\begin{twierdzenie}[Zatrzymany (nad-, pod-)martyngał]
Niech $(M_n)$ będzie martyngałem, podmartyngałem
lub nadmartyngałem względem $(\mathcal F_n)$,
a $\tau$ czasem zatrzymania. Wtedy proces
$M_n^\tau=M_{n\wedge\tau}$ ma tę samą własność
względem $(\mathcal F_n)$.
\end{twierdzenie}
\begin{proof}
Adaptacja wynika z lematu o podstawowych
własnościach zatrzymania. Ponadto
$|M_{n\wedge\tau}|\leq\sum_{k=0}^n|M_k|$,
więc dla każdego $n$ zmienna $M_n^\tau$ jest
całkowalna. Zachodzi tożsamość
\[
M_{(n+1)\wedge\tau}-M_{n\wedge\tau}
=\chi_{\{\tau>n\}}(M_{n+1}-M_n).
\]
Istotnie, na $\{\tau\leq n\}$ obie zatrzymane
wartości są równe $M_\tau$, a na $\{\tau>n\}$
indeksy zatrzymanych wartości to odpowiednio
$n+1$ i $n$, ponieważ czas jest dyskretny.
Zdarzenie $\{\tau>n\}$ należy do $\mathcal F_n$,
więc po warunkowaniu dostajemy
\[
\E[M_{(n+1)\wedge\tau}-M_{n\wedge\tau}
\mid\mathcal F_n]
=\chi_{\{\tau>n\}}\E[M_{n+1}-M_n\mid\mathcal F_n].
\]
Prawa strona jest zerowa dla martyngału,
nieujemna dla podmartyngału i niedodatnia
dla nadmartyngału. Ponieważ $M_{n\wedge\tau}$
jest $\mathcal F_n$-mierzalna, po dodaniu tej
zmiennej do obu stron otrzymujemy odpowiedni
warunek dla procesu zatrzymanego.
\end{proof}

\begin{twierdzenie}[Twierdzenie o próbkowaniu opcjonalnym]
Niech $(M_n)$ będzie martyngałem, a $\sigma,\tau$
czasami zatrzymania spełniającymi
$0\leq\sigma\leq\tau\leq N$ p.n. dla pewnego
ustalonego $N\in\N_0$. Wtedy
$\E[M_\tau\mid\mathcal F_\sigma]=M_\sigma$ p.n.,
a w szczególności $\E M_\tau=\E M_0$.
Dla podmartyngału zachodzi
$\E[M_\tau\mid\mathcal F_\sigma]\geq M_\sigma$
p.n. oraz $\E M_\tau\geq\E M_0$.
Dla nadmartyngału obie nierówności są odwrócone.
\end{twierdzenie}
\begin{proof}
Przyjmujemy wartości umowne na zdarzeniach
nieskończonych czasów zgodnie z definicją.
Z ograniczenia p.n. przez $N$ wynika, że
$|M_\sigma|,|M_\tau|\leq\sum_{k=0}^N|M_k|$
p.n., więc obie zmienne są całkowalne.
Niech $A\in\mathcal F_\sigma$.
Dla każdego $1\leq k\leq N$ funkcja
$H_k=\chi_{A}\chi_{\{\sigma<k\leq\tau\}}$
jest $\mathcal F_{k-1}$-mierzalna:
$A\cap\{\sigma<k\}=A\cap\{\sigma\leq k-1\}
\in\mathcal F_{k-1}$,
a $\{k\leq\tau\}=\{\tau>k-1\}\in\mathcal F_{k-1}$.
Na zdarzeniu $\{\sigma\leq\tau\leq N\}$
pełnej miary zachodzi teleskopowy wzór
\[
\chi_{A}(M_\tau-M_\sigma)
=\sum_{k=1}^N H_k(M_k-M_{k-1}).
\]
Wszystkie składniki są całkowalne, więc
\[
\E[H_k(M_k-M_{k-1})]
=\E[H_k\E[M_k-M_{k-1}\mid\mathcal F_{k-1}]].
\]
Dla martyngału wyrażenie jest zerowe,
a dla pod- i nadmartyngału jest odpowiednio
nieujemne i niedodatnie, ponieważ $H_k\geq0$.
Po zsumowaniu dostajemy dla każdego
$A\in\mathcal F_\sigma$
$\E[\chi_A(M_\tau-M_\sigma)]=0$
dla martyngału, z odpowiednią nierównością
w pozostałych przypadkach. Z mierzalności
$M_\sigma$ względem $\mathcal F_\sigma$
wynika teza dla martyngału na mocy definicji
warunkowej wartości oczekiwanej.
Dla podmartyngału wybieramy w szczególności
$A=\{\E[M_\tau\mid\mathcal F_\sigma]-M_\sigma<0\}$.
Gdyby $\Prob(A)>0$, całka różnicy po $A$
byłaby ujemna, wbrew wykazanej nierówności.
Dla nadmartyngału stosujemy ten argument do $-M$.
Wreszcie stosując udowodniony wynik z $\sigma=0$
i biorąc wartość oczekiwaną, otrzymujemy
odpowiednią relację między $\E M_\tau$ i $\E M_0$.
\end{proof}

\begin{uwaga}[Dlaczego ograniczenie czasu jest ważne]
Powyższe twierdzenie nie mówi, że można zatrzymywać
martyngał w \emph{dowolnym} prawie na pewno
skończonym czasie i zachować jego średnią.
Przejście $N\to\infty$ wymaga kontroli
całkowalności zatrzymanych wartości.
Bez takiej kontroli strategia „czekam aż wygram”
może dawać pozornie sprzeczny wynik.
\end{uwaga}

\begin{twierdzenie}[Zatrzymanie w czasie o skończonej średniej]
Niech $(M_n)$ będzie martyngałem, $\tau$
czasem zatrzymania o $\E\tau<\infty$,
a przyrosty spełniają
$|M_n-M_{n-1}|\leq c$ p.n. dla wszystkich $n\geq1$,
gdzie $0\leq c<\infty$ jest stałą deterministyczną.
Wówczas $M_\tau\in L^1(\Prob)$
oraz $\E M_\tau=\E M_0$.
\end{twierdzenie}
\begin{proof}
Z nieujemności czasu i skończonej średniej
wynika $\Prob(\tau<\infty)=1$.
Ponieważ nierówności dla przyrostów są przeliczalną
rodziną warunków zachodzących p.n., wybieramy
jeden zbiór pełnej miary, na którym zachodzą
wszystkie te nierówności i $\tau<\infty$.
Na tym zbiorze dla każdego $n$
\[
|M_{n\wedge\tau}|
\leq|M_0|+\sum_{k=1}^{n\wedge\tau}|M_k-M_{k-1}|
\leq |M_0|+c\tau.
\]
Na każdej rozważanej trajektorii dla $n\geq\tau$
mamy $M_{n\wedge\tau}=M_\tau$. Stąd
$|M_\tau|\leq|M_0|+c\tau$ p.n., więc $M_\tau$
jest całkowalna. Ponadto
$|M_{n\wedge\tau}-M_\tau|\leq2(|M_0|+c\tau)$
p.n. Prawa strona jest całkowalna, dlatego
twierdzenie o zbieżności zmajoryzowanej daje
$\E|M_{n\wedge\tau}-M_\tau|\to0$. Dla każdego $n$ proces zatrzymany
jest martyngałem i $\E M_{n\wedge\tau}=\E M_0$.
Przechodzimy do granicy w średnich.
\end{proof}


\begin{przyklad}[Zadanie rozwiązane: pierwszy orzeł, ale nie później niż po trzech rzutach]
Niech $\xi_1,\xi_2,\xi_3$ będą niezależnymi uczciwymi
znakami ($+1$ dla orła, $-1$ dla reszki),
$S_n=\sum_{k=1}^n\xi_k$ i
$\tau=\min\{\inf\{n\in\{1,2,3\}:\xi_n=1\},3\}$,
z konwencją $\inf\varnothing=\infty$.
Przyjmujemy $S_0=0$ i
$\mathcal F_n=\sigma(\xi_1,\ldots,\xi_{\min(n,3)})$.
Dla $n<3$ zdarzenie $\{\tau\leq n\}$ oznacza
co najmniej jednego orła w pierwszych $n$ rzutach,
a dla $n\geq3$ jest równe $\Omega$.
Zatem $\tau$ jest czasem zatrzymania i $\tau\leq3$.
Jeśli chcemy traktować $(S_n)$ jako proces
dla wszystkich $n\geq0$, przyjmujemy $S_n=S_3$
dla $n\geq3$. Wszystkie możliwe \emph{zatrzymane} scenariusze
przedstawia tabela:
\begin{center}
\begin{tabular}{c|c|c|c}
zaobserwowane rzuty & $\Prob$ & $\tau$ & $S_\tau$\\\hline
$O$ & $1/2$ & $1$ & $1$\\
$RO$ & $1/4$ & $2$ & $0$\\
$RRO$ & $1/8$ & $3$ & $-1$\\
$RRR$ & $1/8$ & $3$ & $-3$
\end{tabular}
\end{center}
Zatem
\[
\E S_\tau=\frac12+0-\frac18-\frac38=0
=\E S_0,\qquad
\E\tau=\frac12+\frac12+\frac34=\frac74.
\]
Nawet jeśli wybór chwili zatrzymania zależy od
zaobserwowanych wyników, oczekiwany wynik uczciwej
gry nie zmienia się, gdy czas jest ograniczony.
Nie oznacza to, że $S_\tau=0$ na każdej trajektorii.
\end{przyklad}

\begin{przyklad}[Czekanie na jeden punkt]
Niech $S_n=\xi_1+\cdots+\xi_n$ będzie symetrycznym
spacerem losowym na $\Z$ z $S_0=0$,
a $\tau=\inf\{n\geq0:S_n=1\}$.
Z obliczenia ruiny gracza
(rozdział~\ref{ch:markov}), po przesunięciu
przedziału $[-m,1]$ na $\{0,\ldots,m+1\}$
dla całkowitego $m\geq1$, otrzymujemy
$\Prob_0(\tau_1<\tau_{-m})=m/(m+1)$.
Ponieważ zdarzenia po lewej zawierają się
w $\{\tau_1<\infty\}$, dla każdego $m\geq1$
mamy $\Prob(\tau<\infty)\geq m/(m+1)$.
Przejście do granicy daje $\Prob(\tau<\infty)=1$. Jednak $S_\tau=1$
prawie na pewno, więc $\E S_\tau=1\ne0=\E S_0$.
Przyrosty są ograniczone przez $1$;
z poprzedniego twierdzenia wynika, że
$\E\tau=\infty$. To nie jest kontrprzykład
do twierdzenia o zatrzymaniu: zawiodło
właśnie wymagane założenie.
\end{przyklad}



\begin{uwaga}[Dlaczego nie można usunąć ograniczenia czasu bez dowodu]
Niech $S_n$ będzie symetrycznym spacerem losowym i
$\tau=\inf\{n\geq0:S_n=1\}$. W poprzednim przykładzie
wykazano $\tau<\infty$ prawie na pewno, jednak
$\E S_\tau=1\ne0=\E S_0$. Dla
$\tau_m=\tau\wedge m$ mamy na mocy ograniczonego
zatrzymywania $\E S_{\tau_m}=0$ dla każdego $m$.
Jednocześnie $S_{\tau_m}\to1$ prawie na pewno.
Gdyby rodzina $(S_{\tau_m})$ była jednostajnie
całkowalna, zbieżność prawie na pewno i twierdzenie
Vitaliego o zbieżności w $L^1$
(rozdział~\ref{ch:convergence}) dawałyby $\E S_{\tau_m}\to1$, co jest
sprzeczne ze stałą wartością zero.
Zatem rodzina zatrzymanych wartości nie jest
jednostajnie całkowalna. Ogólne twierdzenie
o zatrzymywaniu martyngałów jednostajnie całkowalnych
udowodnimy w rozdziale~\ref{ch:nadmartyngaly}. Przykład wskazuje
\emph{dokładnie}, które przejście graniczne jest
niedozwolone, zamiast poprzestawać na stwierdzeniu,
że $\tau$ jest nieograniczony.
\end{uwaga}


Podstawowy rachunek martyngałowy, transformacje
przewidywalne i nierówności maksymalne rozwiniemy
w rozdziale~\ref{ch:martyngaly}; kryteria zbieżności
i dalsze wersje twierdzeń o zatrzymywaniu
w rozdziale~\ref{ch:nadmartyngaly}.
